\scanpage{150}
\[
\begin{aligned}
g\cdot\tau=\tau\circ\tau_\infty(g^{-1})\qquad \text{où}\quad
&\tau_\infty=\begin{pmatrix}-1&0\\0&1\end{pmatrix}\in\mathrm{Gl}(2,\mathbb Z)\\
&\tau_\infty(g^{-1})=\begin{pmatrix}d&b\\c&a\end{pmatrix}\quad \text{si}\\
&\hspace{5em}g=\begin{pmatrix}a&b\\c&d\end{pmatrix}
\end{aligned}
\]

\uline{Corollaire 1} Le groupe fondamental $\unsure{\mathcal T|\mathscr T}_{1,1}$ de $M_{11}$ relatif au revêtement universel précédent est canoniquement isomorphe à $SL(2,\mathbb Z)$ (avec l'opération précédente \inserted{à gauche} de $SL(2,\mathbb Z)$ sur ledit revêtement, déduite de l'opération canonique \inserted{\unsure{à|d}} droite par le passage de $g$ à $g^{-1}$ .-- ).

Les marteaux-pilons sur les $\mathcal T_{g,\unsure{\nu|v}}$ donnent alors:

\uline{Corollaire 2} Soit $(\deleted{\unsure{\Sigma|Z|E}}X,o)$ une surface orientée compacte de genre 1, $\mathcal T$ le groupe des classes d'isotopie d'automorphismes de $(X,o)$, $\deleted{\unsure{d|o}}H=\pi_1(X,o)=H_1(X,\mathbb Z)$, \deleted{\unsure{considérons}}\inserted{alors} \unsure{l'|le} homomorphisme canonique
\[
\mathcal T\longrightarrow\mathrm{Aut}^{+}_{\mathbb Z}(H)\tag*{(16)}
\]
est un isomorphisme

\uline{Corollaire 3} Soient $\pi$ un groupe à bouts du type $(1,1)$, $H=\pi_{\mathrm{ab}}$, $\mathcal T$ le groupe des autom. ext à bouts de $\pi$ compatibles avec l'orientation, alors l'hom canonique
\[
\mathcal T\longrightarrow\mathrm{Aut}^{+}_{\mathbb Z}(H)\qquad\text{est un iso.}
\]
% archive sheet number: 149
